\subsection{The limit of the exponential equations}
\label{subsec:semipositive-exponential-setting}

Use the data of \cref{thm:semipositive-dual-exact-equation}: the manifold $X$ of dimension $n\geq2$, Hermitian form $\omega$, probability volume $\mathrm{d}V$, nonzero nonnegative density $f$, positive threshold $b=b_\omega(f)$, and smooth semipositive maximizing form $T$ with analytic degeneracy set $Z$. Put $p=f^{1/n}$ and let $k_T$ be its determinant density. All $L^1$ convergence below uses $\mathrm{d}V$.

A smooth semipositive maximizing form also identifies the normalized exponential limit. The argument passes the linear trace inequalities to the limit before using the local regularity proof.

\begin{proposition}\label{prop:semipositive-dual-exponential-limit}
Assume the hypotheses of \cref{thm:semipositive-dual-exact-equation}, and let $u_0$ be its unique threshold solution with $\sup_Xu_0=0$. For real $\lambda>0$, let $u_\lambda$ be the greatest exponential solution of \cref{cor:individual-positive-threshold-exponential}, and put
\[
 v_\lambda=u_\lambda-\sup_Xu_\lambda,\qquad
 c_\lambda=\mathrm{e}^{\lambda\sup_Xu_\lambda}.
\]
There is a constant $C$ independent of $\lambda>0$ such that
\[
 b\leq c_\lambda\leq b\,\mathrm{e}^{C\lambda},\qquad
 u_0-C\leq v_\lambda\leq0.
\]
Moreover $v_\lambda\to u_0$ in $L^1(X)$ as $\lambda\downarrow0$.
\end{proposition}
\begin{proof}
Write $p=f^{1/n}$ and introduce the probability measure
\[
 \mathrm{d}\nu=b^{1/n}p k_T\,\mathrm{d}V.
\]
It has mass one because $T$ maximizes. Its density belongs to $L^n$, since $p\in L^n$ and the smooth semipositive form $T$ has bounded determinant density $k_T$. By \cref{lem:uniform-exponential}, normalized $\omega$-plurisubharmonic functions have uniformly bounded moments of every finite order. H\"older's inequality with exponents $n$ and $n/(n-1)$ therefore gives
\[
 \int_X(-v_\lambda)\,\mathrm{d}\nu\leq C.
\]
Here $n\geq2$ is used.

Apply the trace inequality to $u_\lambda$, or equivalently to $v_\lambda$, and pass from positive regularizations of $T$ to $T$ as in the preceding proof. After integration,
\[
 1\geq(c_\lambda/b)^{1/n}
       \int_X\mathrm{e}^{\lambda v_\lambda/n}\,\mathrm{d}\nu
 \geq(c_\lambda/b)^{1/n}\mathrm{e}^{-C\lambda/n},
\]
where the last step is Jensen's inequality. This proves the upper scalar bound. The lower bound and the comparison
\[
 v_\lambda\geq u_0-\lambda^{-1}\log(c_\lambda/b)
\]
are \cref{cor:individual-positive-threshold-exponential}, giving the asserted common lower barrier.

Take any sequence $\lambda_j\downarrow0$ and an $L^1$-convergent subsequence $v_{\lambda_j}\to v$, supplied by \cref{thm:normalized-compactness}. Its limit has supremum zero. The normalized $L^1$ bound implies $\lambda v_\lambda\to0$ in measure. Since $v_\lambda\leq0$ and $p\in L^1$, it follows that
\[
 c_\lambda^{1/n}\mathrm{e}^{\lambda v_\lambda/n}p
       \longrightarrow b^{1/n}p
       \quad\text{in }L^1(\mathrm{d}V).
\]
Put $A_\lambda=\omega+\ddc v_\lambda$. The positive trace errors
\[
 A_\lambda\wedge T
   -c_\lambda^{1/n}\mathrm{e}^{\lambda v_\lambda/n}p k_T\,\mathrm{d}V
\]
have masses tending to zero, since the two terms have total masses tending to one. The $L^1$ convergence of the potentials gives weak convergence of their linear currents. Consequently
\[
 (\omega+\ddc v)\wedge T=b^{1/n}p k_T\,\mathrm{d}V.
\]
For every smooth Hermitian form $\beta$, the all-Hermitian trace inequality passes to the same limit and gives
\[
 (\omega+\ddc v)\wedge\beta^{n-1}
   \geq(bf\,\mathrm{d}V/\beta^n)^{1/n}\beta^n.
\]
The local matrix and regularity argument in the proof of \cref{thm:semipositive-dual-exact-equation} now applies on $X\setminus Z$: with $\gamma$ the positive root of $T$ there and $a=(bf\,\mathrm{d}V/\gamma^n)^{1/n}$, it gives $\omega+\ddc v=a\gamma$, then smoothness, and then the exact equation there. Both extended measures ignore $Z$, so $v$ is a global threshold solution. Uniqueness and normalization give $v=u_0$.

Every sequence tending to zero has a subsequence with this same $L^1$ limit. Normalized compactness therefore proves convergence of the whole family.
\end{proof}

The lower barrier gives capacity convergence under the additional envelope condition below. If the maximizing form is strictly positive everywhere, \cref{cor:smooth-dual-exponential-limit} already gives uniform convergence.
